\documentclass[11pt,a4paper]{article}
\usepackage[margin=2.5cm]{geometry}
\usepackage{fontspec}
\usepackage[brazil]{babel}
\usepackage{booktabs}
\usepackage{hyperref}
\hypersetup{colorlinks=true,linkcolor=black}
\setlength{\parindent}{1.25cm}
\setlength{\parskip}{0.25em}

\title{MAC0219 - Programação Concorrente e Paralela\\Mini EP4 -- Profiling e otimização}
\author{Leonardo Heidi Almeida Murakami\\NUSP: 11260186}
\date{26 de agosto de 2026}

\begin{document}
\maketitle

\section{Profiling antes das otimizações}

O programa original foi analisado com \texttt{gprof} e \texttt{perf}. As duas ferramentas mostraram que praticamente todo o tempo era gasto nas funções \texttt{slowsort} e \texttt{fibonacci}.

\begin{table}[ht]
\centering
\begin{tabular}{lrr}
\toprule
Função & gprof & perf \\
\midrule
\texttt{slowsort} & 57,55\% & 49,70\% \\
\texttt{fibonacci} & 42,45\% & 50,27\% \\
\bottomrule
\end{tabular}
\end{table}

O problema era a quantidade de chamadas recursivas. O \texttt{gprof} registrou cerca de 342 milhões de chamadas recursivas de \texttt{slowsort} e 459 milhões de \texttt{fibonacci}. Isso acontece porque o Slowsort é um algoritmo de ordenação muito lento e porque os mesmos números de Fibonacci eram calculados várias vezes.

\section{Otimizações}

Foram feitas duas mudanças:

\begin{itemize}
  \item O \texttt{slowsort} foi trocado pelo \texttt{qsort}, que é muito mais eficiente e evita milhões de chamadas recursivas.
  \item Os números de Fibonacci de 0 a 35 passaram a ser calculados uma única vez e guardados em um vetor. Depois disso, cada valor pode ser consultado diretamente, sem repetir os cálculos.
\end{itemize}

Essas mudanças atacam exatamente as duas funções que apareceram como gargalos no profiling. As duas versões continuaram produzindo o mesmo resultado:

\begin{center}
\texttt{Soma dos números de Fibonacci pares: 20100614}
\end{center}

\newpage

\section{Resultados depois das otimizações}

A versão original levou, em média, 0,943 s no \texttt{perf stat}. A versão otimizada levou 0,527 ms. Isso representa uma aceleração de aproximadamente $1.790\times$ e uma redução de 99,94\% no tempo de execução.

\begin{table}[ht]
\centering
\begin{tabular}{lrr}
\toprule
Métrica & Original & Otimizado \\
\midrule
Tempo & 0,943 s & 0,527 ms \\
Ciclos & 4.723.080.682 & 272.605 \\
Instruções & 17.030.478.110 & 228.259 \\
Desvios & 2.859.087.265 & 43.253 \\
\bottomrule
\end{tabular}
\end{table}

O \texttt{gprof} não conseguiu acumular tempo na versão otimizada, pois ela terminou antes do intervalo de amostragem de 0,01 s. O \texttt{perf record} coletou apenas 11 amostras, concentradas principalmente na inicialização e no encerramento do programa. Portanto, depois das mudanças, nenhuma função do código apareceu como um novo gargalo.

\end{document}
